\documentclass{article}
\usepackage{graphicx}
\usepackage[margin=2cm, a4paper]{geometry}
\usepackage{setspace}
\usepackage{pdfpages}
\usepackage{float}
\usepackage{amsmath}
\usepackage{fancyvrb}
\usepackage{amssymb}
\usepackage{minted}
\usepackage{enumitem}
\usepackage[colorlinks,linkcolor=blue]{hyperref}


\renewcommand{\contentsname}{Contents}
\renewcommand{\figurename}{Figure}
\renewcommand{\tablename}{Table}
\hypersetup{
    colorlinks=true,
    linkcolor=black,
    filecolor=magenta,      
    urlcolor=blue,
}
\newcommand{\mytitle}{CSIE 5452, Spring 2026: Homework 1}
\newcommand{\myauthor}{B13902022 Yu-Chi,Lai}

\usepackage{fancyhdr}
\pagestyle{fancy}

\fancyhead{}
\fancyhead[L]{\mytitle}
\fancyhead[R]{\myauthor}

\title{\mytitle}
\author{\textbf{\myauthor}}
\date{}


\begin{document}

\maketitle

\section{Timing Analysis of the CAN Protocol: Part I (12pt)}
\subsection{}
Firstly, set $Q_0=B_0=C_1=30$, since no other task has higher priority than $\mu_0$, the terms in the summation is zero.

Hence, the worst-case response time $R_0=Q_0+C_0=30+10=40$(msec).
\begin{table}[h]
    \centering % This centers the table on the page
    \begin{tabular}{|c||c||c|c|c|} 
        \hline 
        Iteration & LHS ($Q_0$) & $B_0$ & RHS & Stop? \\ [0.5ex] 
        \hline \hline
        1 & 30 & 30 & 30 & yes\\ [1ex] 
        \hline 
    \end{tabular}
\end{table}

\subsection{}
Let $Q_1=B_1=C_1=30$ as the beginnning.
\begin{table}[h]
    \centering
    % Adding extra row height to prevent the fraction from touching the lines
    \renewcommand{\arraystretch}{1.5} 
    
    \begin{tabular}{|c||c||c|c|c|c|c|c|c|c|}
        \hline
        Iteration & LHS ($Q_1$) & $B_1$ & $j$ & $Q_1 + \tau$ & $T_j$ & $\left\lceil \frac{Q_1 + \tau}{T_j} \right\rceil$ & $C_j$ & RHS & Stop? \\ 
        \hline \hline
        1 & 30 & 30 & 0 & 30.1 & 50 & 1 & 10 & 40 & no \\ 
        \hline
        2 & 40 & 30 & 0 & 40.1 & 50 & 1 & 10 & 40 & yes \\ 
        \hline
    \end{tabular}
\end{table}

Thus, the worst-case response time of $\mu_1$ is $R_1=Q_1+C_1=40+30=70$(msec).

\subsection{}
Let $Q_2=B_2=C_2=20$ as the beginnning.
\begin{table}[h]
    \centering
    % Adding extra row height to prevent the fraction from touching the lines
    \renewcommand{\arraystretch}{1.5} 
    
    \begin{tabular}{|c||c||c|c|c|c|c|c|c|c|}
        \hline
        Iteration & LHS ($Q_2$) & $\left\lceil \frac{Q_2 + \tau}{T_0} \right\rceil C_0+\left\lceil \frac{Q_2+\tau}{T_1}\right\rceil C_1$ & RHS & Stop? \\ 
        \hline \hline
        1 & 20 & 40 & 60 & no \\ 
        \hline
        2 & 60 & 50  & 70 & no \\ 
        \hline
        3 & 70 & 50 & 70 & yes \\ 
        \hline
    \end{tabular}
\end{table}
After 3 iterations, $Q_2$ converges. Hence, the worst case response time $R_2=Q_2+C_2=70+20=90$(msec).
\section{Timing Analysis of the CAN Protocol: Part II (36pt)}

\subsection{}
\begin{align*}
R_{0} &= 1 \\
R_{1} &= 2 \\
R_{2} &= 2 \\
R_{3} &= 3 \\
R_{4} &= 3 \\
R_{5} &= 4 \\
R_{6} &= 5  \\
R_{7} &= 8 \\
R_{8} &= 9 \\
R_{9} &= 9 \\
R_{10} &= 10 \\
R_{11} &= 19 \\
R_{12} &= 19 \\
R_{13} &= 20 \\
R_{14} &= 29 \\
R_{15} &= 29 \\
R_{16} &= 30
\end{align*}
\subsection{}
\begin{minted}[frame=lines,framesep=2mm,baselinestretch=1.2,linenos,breaklines]{cpp}
#include <bits/stdc++.h>
using namespace std;
signed main() {
    ifstream fin("input.dat");
    int n;
    double tau;
    fin >> n;
    fin >> tau;
    vector<int> priority(n);
    vector<double> transmission_time(n), period(n);
    for (int i = 0; i < n; ++i) {
        fin >> priority[i] >> transmission_time[i] >> period[i];
    }
    fin.close();
    vector<int> res(n);
    for (int i = 0; i < n; i++) {
        double block = 0;
        for (int j = 0; j < n; j++) {
            if (priority[j] >= priority[i]) {
                block = max(block, transmission_time[j]);
            }
        }
        double wait = block;
        while (1) {
            double result  = block;
            for (int j = 0; j < n; j++) {
                if (priority[j] < priority[i]) {
                    result += transmission_time[j] * ceil((wait + tau) / period[j]); 
                }
            }
            if (result > period[i]) {
                res[i] = -1;
                break;
            }
            if (result-wait <= 1e-9 && result-wait >= -(1e-9)) {
                res[i] = transmission_time[i] + wait;
                break;
            } else {
                wait = result;
                continue;
            }
        }
    }
    for (int i = 0; i < n; i++) {
        cout << res[i] << endl;
    }
}
\end{minted}

\section{Timing Analysis of Preemptive Fixed-Priority Scheduling (16pt)}
\section{Timing Analysis of TDMA-Based Protocols (12pt)}

\section{MILP Linearization (12pt)}
\section{Signal Packing (12pt)}


\end{document}
% how to display codes?
% \begin{minted}[frame=lines,framesep=2mm,baselinestretch=1.2,linenos,breaklines]{python}
% \end{minted}

% how to display images?
% \begin{figure}[H]
%     \centering
%     \includegraphics[width=0.5\linewidth]{}
%     \caption{Caption}
% \end{figure}
% test